\documentclass{article}
\usepackage[preprint]{neurips_2023}

\usepackage[utf8]{inputenc} % allow utf-8 input
\usepackage[T1]{fontenc}    % use 8-bit T1 fonts
\usepackage{hyperref}       % hyperlinks
\usepackage{url}            % simple URL typesetting
\usepackage{booktabs}       % professional-quality tables
\usepackage{amsfonts}       % blackboard math symbols
\usepackage{nicefrac}       % compact symbols for 1/2, etc.
\usepackage{microtype}      % microtypography
\usepackage{xcolor}         % colors

\usepackage{amsmath}
\usepackage{graphicx}
\usepackage{algorithm}
\usepackage[noend]{algpseudocode}

\DeclareMathOperator{\softmax}{softmax}
\DeclareMathOperator{\sigmoid}{sigmoid}
\DeclareMathOperator{\sign}{sign}
\DeclareMathOperator{\relu}{ReLU}

\title{Deep Convolutional Neural Networks for Semantic Segmentation}

\author{
  Merrick Qiu\\
  UC San Diego\\
  La Jolla, CA 92093 \\
  \texttt{myqiu@ucsd.edu} \\
  \And
  Michael Ye \\
  UC San Diego \\
  La Jolla, CA 92093 \\
  \texttt{mhye@ucsd.edu} \\
  \AND
  Collin Cranston \\
  UC San Diego \\
  La Jolla, CA 92093 \\
  \texttt{YOUR-EMAIL@ucsd.edu} \\
}



\begin{document}


\maketitle


\begin{abstract}
  We used a deep convolutional neural network in order to perform semantic segmentation on
  the PASCAL VOC 2012 dataset. 
  We built a PyTorch CNN with ReLU activation function, batch normalization, and an encoder and decoder architecture.
  With the baseline model we were able to get a mean pixel accuracy of [PIXEL ACCURACY] and mean IoU score of [IOU] but after using [SET OF OPTIMIZATIONS] 
  we were able to get a model with a mean pixel accuracy of [PIXEL ACCURACY] and a mean IoU score of [IOU].
  It was interesting that [FINDINGS].
\end{abstract}


\section{Introduction}
\subsection{Overview}
Semantic segmentation is the process by which the pixels in an image are classified into a number of categories.
In the PASCAL VOC 2012 dataset, each pixel is classified into one of 21 classes using a softmax.
This is useful in a number of fields including self-driving cars, medical imaging, and satellite imaging.
However, this is a challenging task because it requires the model to understand the context of the image and to be able to
distinguish between different objects. 

To perform semantic segmentation, we used a deep convolutional neural network with an encoder and decoder architecture.
Convolutions can calculate the features of an image by repeatedly applying a filter to an image.
We use ReLU as our acitvation function to add nonlinearity to the network. 
Batch normalization normalizes the input to each layer so that it has zero mean and scaling it to the standard deviation.
This speeds up training and makes the model more robust.
To perform a feedforward pass,
the encoder is used to extract features from the image and the decoder is used to upsample the features to the original image size.


\subsection{Xavier and Kaiming Weight Initialization}
The output from a matrix multiplication has a variance equal to the 
the size of the input vector if the weights and inputs 
are all initated with standard normal distributions.
In addition, the activation function can change the variance of the output.
In order to prevent activations from exploding or vanishing during the forward pass of the initial model,
the weights are set using Kaiming weight initialization since we are using ReLU.
Xavier weight initialization is used when the activation function is a sigmoid or hyperbolic tangent.

We want the variance of the layer outputs to be $1$ so the matrix weights
are initialized with a standard deviation of $\sqrt{1/n}$.
However the ReLU activation function halves the variance of the output
so instead we want to initialize the weights with standard deviation of $\sqrt{2/n}$.
Although in the original paper by Kaiming He et al.
they used the normal distribution $\mathcal{N}(0, \sqrt{2/n})$,
\href{https://pytorch.org/docs/stable/nn.init.html\#torch.nn.init.kaiming\_uniform\_}{PyTorch}
uses the uniform distribution $\mathcal{U}(-\sqrt{6/n}, \sqrt{6/n})$.

\subsection{CNN architecture}

The architecture of the basline model is dhown in table 1.
Note that the output is passed through ReLU after each convolution or transpose convolution.
\begin{table}[h!]
\caption{Baseline model architecture}
\label{tab:baseline_model_architecture}
\centering
\begin{tabular}{ccc}
  \toprule
  \textbf{Type} & \textbf{Output Channels} & \textbf{Kernel Size / Stride / Padding} \\ 
  \midrule
  Conv2d & 32 & 3x3 / 2 / 1 \\ 
  BatchNorm2d & 32 & - \\ 
  Conv2d & 64 & 3x3 / 2 / 1 \\ 
  BatchNorm2d & 64 & - \\ 
  Conv2d & 128 & 3x3 / 2 / 1 \\ 
  BatchNorm2d & 128 & - \\ 
  Conv2d & 256 & 3x3 / 2 / 1 \\ 
  BatchNorm2d & 256 & - \\ 
  Conv2d & 512 & 3x3 / 2 / 1 \\ 
  BatchNorm2d & 512 & - \\ 
  ConvTranspose2d & 512 & 3x3 / 2 / 1 / 1 \\ 
  BatchNorm2d & 512 & - \\ 
  ConvTranspose2d & 256 & 3x3 / 2 / 1 / 1 \\ 
  BatchNorm2d & 256 & - \\ 
  ConvTranspose2d & 128 & 3x3 / 2 / 1 / 1 \\ 
  BatchNorm2d & 128 & - \\ 
  ConvTranspose2d & 64 & 3x3 / 2 / 1 / 1 \\ 
  BatchNorm2d & 64 & - \\ 
  ConvTranspose2d & 32 & 3x3 / 2 / 1 / 1 \\ 
  BatchNorm2d & 32 & - \\ 
  Conv2d & n\_class & 1x1 / - / - \\ 
  \bottomrule
\end{tabular}
\end{table}
\section{Related Work}

\section*{References}
% He paper
% https://arxiv.org/pdf/1502.01852


\end{document}